\documentclass[12pt,a4paper]{article}
\usepackage[utf8]{inputenc}
\usepackage[T5]{fontenc}
\usepackage{amsmath, amssymb}
\usepackage{graphicx}
\usepackage{ifthen}

\newboolean{showsolutions}
\setboolean{showsolutions}{true}

% --- ĐỊNH NGHĨA CÁC LỆNH ẢO CHO CÔNG CỤ LATEX2JSON CỦA WEBSITE ---
\newcommand{\doctitle}[1]{\begin{center}\Large\textbf{#1}\end{center}}
\newcommand{\docdesc}[1]{\textbf{Mô tả:} #1}
\newcommand{\docgrade}[1]{\textbf{Lớp:} #1}
\newcommand{\docstatus}[1]{\textbf{Trạng thái:} #1}
\newcommand{\doctype}[1]{\textbf{Loại:} #1}
\newcommand{\doctopics}[1]{\textbf{Chủ đề:} #1}

\newenvironment{textblock}{\vspace{1em}}{}

\newenvironment{quiz}[1]{
	\vspace{1em}\noindent\textbf{\Large Kiểm tra: #1}
}{}

\newenvironment{mcq}[4]{
	\vspace{1em}\noindent\textbf{Câu hỏi:} #1 \\
	\ifthenelse{\boolean{showsolutions}}{
		\textit{(Đáp án đúng: #2, Điểm: #3) - Giải thích: #4}
	}{}
	\begin{itemize}
	}{
	\end{itemize}
}
\newcommand{\option}[1]{\item #1}

\newenvironment{truefalse}[3]{
	\vspace{1em}\noindent\textbf{Đúng/Sai:} #1 \\
	\ifthenelse{\boolean{showsolutions}}{
		\textit{(Điểm: #2) - Giải thích: #3}
	}{}
	\begin{itemize}
	}{
	\end{itemize}
}
\newcommand{\statement}[2]{\item [\textbf{#1}] #2}

\newcommand{\shortanswer}[4]{
    \vspace{1em}\noindent\textbf{Trả lời ngắn (Điểm: #3):}

    \noindent #1

	\ifthenelse{\boolean{showsolutions}}{
		\vspace{0.5em}
		\noindent\textit{Đáp án: #2 - Giải thích:}
		\noindent #4
	}{}
}

\begin{document}

\doctitle{PHẦN 03: BÀI TẬP RÈN LUYỆN - PHƯƠNG TRÌNH ĐƯỜNG THẲNG}
\docdesc{Bộ câu hỏi trắc nghiệm nhiều phương án lựa chọn, trắc nghiệm đúng/sai và trả lời ngắn chuyên đề Phương trình đường thẳng trong không gian Oxyz}
\docgrade{12}
\docstatus{draft}
\doctype{quiz}
\doctopics{Hình học không gian, Oxyz, Phương trình đường thẳng, Trắc nghiệm rèn luyện}

% =========================================================================
% 1. CÂU TRẮC NGHIỆM NHIỀU PHƯƠNG ÁN LỰA CHỌN
% =========================================================================

\begin{quiz}{Phần 1. Trắc nghiệm nhiều phương án lựa chọn}

\begin{mcq}{Trong không gian $Oxyz$, cho đường thẳng $\Delta: \begin{cases} x = 2 + 3t \\ y = -1 + 6t \\ z = 5 - 9t \end{cases}$ ($t \in \mathbb{R}$). Vecto nào sau đây là một vecto chỉ phương của đường thẳng $\Delta$?}{2}{1}{Đường thẳng $\Delta$ có một VTCP $\vec{u} = (3; 6; -9) = 3(1; 2; -3)$. Vậy $\vec{u_3} = (1; 2; -3)$ là một VTCP của $\Delta$. Chọn C.}
	\option{$\vec{u_1} = (2; -1; 5)$.}
	\option{$\vec{u_2} = (-1; 2; 3)$.}
	\option{$\vec{u_3} = (1; 2; -3)$.}
	\option{$\vec{u_4} = (1; -2; 3)$.}
\end{mcq}

\begin{mcq}{Trong không gian $Oxyz$, cho đường thẳng $d: \frac{x}{4} = \frac{y-1}{-2} = \frac{z+1}{1}$. Một vecto chỉ phương của đường thẳng $d$ có tọa độ là}{3}{1}{Đường thẳng $d$ có một VTCP $\vec{u} = (4; -2; 1)$. Chọn D.}
	\option{$(0; 1; -1)$.}
	\option{$(4; 2; 1)$.}
	\option{$(0; -1; 1)$.}
	\option{$(4; -2; 1)$.}
\end{mcq}

\begin{mcq}{Trong không gian $Oxyz$, cho đường thẳng $d: \begin{cases} x = 1 \\ y = 2 + 3t \\ z = 5 - t \end{cases}$ ($t \in \mathbb{R}$). Điểm nào dưới đây thuộc đường thẳng $d$?}{2}{1}{Với $t = 0$, ta có điểm $P(1; 2; 5)$ thuộc đường thẳng $d$. Chọn C.}
	\option{$M(1; 3; -1)$.}
	\option{$N(2; 3; -1)$.}
	\option{$P(1; 2; 5)$.}
	\option{$Q(1; 5; 4)$.}
\end{mcq}

\begin{mcq}{Trong không gian $Oxyz$, cho đường thẳng $\Delta: \frac{x-1}{1} = \frac{y+2}{-1} = \frac{z}{2}$. Điểm nào sau đây thuộc đường thẳng $\Delta$?}{2}{1}{Đường thẳng $\Delta$ đi qua điểm $A(1; -2; 0)$. Chọn C.}
	\option{$M(-1; 2; 0)$.}
	\option{$N(1; -1; 2)$.}
	\option{$A(1; -2; 0)$.}
	\option{$B(1; 2; 0)$.}
\end{mcq}

\begin{mcq}{Trong không gian $Oxyz$, cho đường thẳng $d: \begin{cases} x = -3 + 2t \\ y = 1 - t \\ z = -1 + 4t \end{cases}$ ($t \in \mathbb{R}$). Điểm nào sau đây \textbf{không} thuộc đường thẳng $d$?}{2}{1}{Thay tọa độ điểm $P(-5; 2; -5)$ vào ta được $\begin{cases} -5 = -3 + 2t \\ 2 = 1 - t \\ -5 = -1 + 4t \end{cases} \Leftrightarrow \begin{cases} t = -1 \\ t = -1 \\ t = -1 \end{cases}$ thỏa mãn, tuy nhiên xét các điểm khác: với $t = 1 \Rightarrow (-1; 0; 3)$, với $t = -1 \Rightarrow (-5; 2; -5)$, với $t = 2 \Rightarrow (1; -1; 7)$. Điểm không thuộc là điểm có tọa độ không thỏa mãn phương trình. Chọn C.}
	\option{$M(-3; 1; -1)$.}
	\option{$N(-1; 0; 3)$.}
	\option{$P(-5; 2; -5)$.}
	\option{$Q(1; -1; 7)$.}
\end{mcq}

\begin{mcq}{Trong không gian $Oxyz$, cho đường thẳng $\Delta: \frac{x+4}{3} = \frac{y+2}{2} = \frac{z-4}{-1}$. Điểm nào sau đây \textbf{không} thuộc đường thẳng $\Delta$?}{3}{1}{Thay tọa độ điểm $G(-7; -4; 5)$ vào ta được $\frac{-7+4}{3} = -1; \frac{-4+2}{2} = -1; \frac{5-4}{-1} = -1 \ne 1$ nên $G \notin \Delta$. Chọn D.}
	\option{$E(-4; -2; 4)$.}
	\option{$F(-1; 0; 3)$.}
	\option{$H(2; 2; 2)$.}
	\option{$G(-7; -4; 5)$.}
\end{mcq}

\begin{mcq}{Trong không gian $Oxyz$, đường thẳng đi qua điểm $A(3; 2; 5)$ và nhận $\vec{u} = (-2; 8; -7)$ làm vecto chỉ phương có phương trình tham số là}{2}{1}{Phương trình tham số là $\begin{cases} x = 3 - 2t \\ y = 2 + 8t \\ z = 5 - 7t \end{cases}$ ($t \in \mathbb{R}$). Chọn C.}
	\option{$\begin{cases} x = -2 + 3t \\ y = 8 + 2t \\ z = -7 + 5t \end{cases}$.}
	\option{$\begin{cases} x = 3 - 2t \\ y = 2 + 8t \\ z = -5 - 7t \end{cases}$.}
	\option{$\begin{cases} x = 3 - 2t \\ y = 2 + 8t \\ z = 5 - 7t \end{cases}$.}
	\option{$\begin{cases} x = 3 + 2t \\ y = 2 + 8t \\ z = 5 - 7t \end{cases}$.}
\end{mcq}

\begin{mcq}{Trong không gian $Oxyz$, phương trình chính tắc của đường thẳng đi qua điểm $B(-1; 3; 6)$ và nhận $\vec{u} = (2; -3; 8)$ làm vecto chỉ phương là}{1}{1}{Phương trình chính tắc là $\frac{x+1}{2} = \frac{y-3}{-3} = \frac{z-6}{8}$. Chọn B.}
	\option{$\frac{x-1}{2} = \frac{y+3}{-3} = \frac{z+6}{8}$.}
	\option{$\frac{x+1}{2} = \frac{y-3}{-3} = \frac{z-6}{8}$.}
	\option{$\frac{x+1}{2} = \frac{y-3}{3} = \frac{z-6}{8}$.}
	\option{$\frac{x-2}{-1} = \frac{y+3}{3} = \frac{z-8}{6}$.}
\end{mcq}

\begin{mcq}{Trong không gian $Oxyz$, cho hai điểm $M(1; -1; 1)$ và $N(5; 5; 1)$. Đường thẳng $MN$ có phương trình tham số là}{2}{1}{Ta có $\vec{MN} = (4; 6; 0) = 2(2; 3; 0)$. Đường thẳng qua $M(1; -1; 1)$ nhận $\vec{u} = (2; 3; 0)$ làm VTCP có phương trình là $\begin{cases} x = 1 + 2t \\ y = -1 + 3t \\ z = 1 \end{cases}$. Chọn C.}
	\option{$\begin{cases} x = 1 + 2t \\ y = -1 + 3t \\ z = 1 + t \end{cases}$.}
	\option{$\begin{cases} x = 5 + 2t \\ y = 5 + 3t \\ z = 1 + t \end{cases}$.}
	\option{$\begin{cases} x = 1 + 2t \\ y = -1 + 3t \\ z = 1 \end{cases}$.}
	\option{$\begin{cases} x = 1 + 4t \\ y = -1 + 6t \\ z = 1 + t \end{cases}$.}
\end{mcq}

\begin{mcq}{Trong không gian $Oxyz$, cho hai điểm $A(1; -2; -3), B(-1; 4; 1)$ và đường thẳng $d: \frac{x+2}{1} = \frac{y-2}{-1} = \frac{z+3}{2}$. Phương trình chính tắc của đường thẳng đi qua trung điểm của đoạn thẳng $AB$ và song song với đường thẳng $d$ là}{3}{1}{Trung điểm của $AB$ là $I(0; 1; -1)$. Đường thẳng qua $I$ song song với $d$ có VTCP $\vec{u} = (1; -1; 2)$ nên có phương trình là $\frac{x}{1} = \frac{y-1}{-1} = \frac{z+1}{2}$. Chọn D.}
	\option{$\frac{x-1}{1} = \frac{y+1}{-1} = \frac{z-1}{2}$.}
	\option{$\frac{x}{1} = \frac{y-1}{1} = \frac{z+1}{2}$.}
	\option{$\frac{x}{1} = \frac{y+1}{-1} = \frac{z-1}{2}$.}
	\option{$\frac{x}{1} = \frac{y-1}{-1} = \frac{z+1}{2}$.}
\end{mcq}

\begin{mcq}{Trong không gian $Oxyz$, đường thẳng $\Delta$ đi qua điểm $M(2; 0; -1)$ và vuông góc với mặt phẳng $(P): 2x - 3y + z + 1 = 0$ có phương trình tham số là}{2}{1}{VTCP của $\Delta$ là $\vec{u} = \vec{n}_{(P)} = (2; -3; 1)$. Phương trình tham số là $\begin{cases} x = 2 + 2t \\ y = -3t \\ z = -1 + t \end{cases}$. Chọn C.}
	\option{$\begin{cases} x = 2 + 2t \\ y = -3 \\ z = -1 + t \end{cases}$.}
	\option{$\begin{cases} x = -2 + 2t \\ y = -3t \\ z = 1 + t \end{cases}$.}
	\option{$\begin{cases} x = 2 + 2t \\ y = -3t \\ z = -1 + t \end{cases}$.}
	\option{$\begin{cases} x = 2 + 2t \\ y = 3t \\ z = -1 + t \end{cases}$.}
\end{mcq}

\begin{mcq}{Trong không gian $Oxyz$, đường thẳng $d: \frac{x}{2} = \frac{y-2}{1} = \frac{z+3}{3}$ vuông góc với mặt phẳng nào dưới đây?}{3}{1}{Đường thẳng $d$ có VTCP $\vec{u} = (2; 1; 3)$. Mặt phẳng $(\alpha_4): 2x + y + 3z - 2024 = 0$ có VTPT $\vec{n} = (2; 1; 3) = \vec{u}$ nên $d \perp (\alpha_4)$. Chọn D.}
	\option{$(\alpha_1): 2x + y - 3z + 1 = 0$.}
	\option{$(\alpha_2): x + 2y + 3z - 1 = 0$.}
	\option{$(\alpha_3): 2x - y + 3z + 2 = 0$.}
	\option{$(\alpha_4): 2x + y + 3z - 2024 = 0$.}
\end{mcq}

\begin{mcq}{Trong không gian $Oxyz$, cho hai đường thẳng $d_1: \begin{cases} x = 1 + 2t \\ y = 2 + 3t \\ z = 3 + 4t \end{cases}$ ($t \in \mathbb{R}$) và $d_2: \begin{cases} x = 3 + 4t' \\ y = 5 + 6t' \\ z = 7 + 8t' \end{cases}$ ($t' \in \mathbb{R})$. Mệnh đề nào sau đây đúng?}{2}{1}{$\vec{u_1} = (2; 3; 4), \vec{u_2} = (4; 6; 8) = 2\vec{u_1}$. Điểm $A(1; 2; 3) \in d_1$ thay vào $d_2$ ta được $t' = -\frac{1}{2}$ (thỏa mãn). Vậy $d_1 \equiv d_2$. Chọn C.}
	\option{$d_1 \perp d_2$.}
	\option{$d_1 \parallel d_2$.}
	\option{$d_1 \equiv d_2$.}
	\option{$d_1$ và $d_2$ chéo nhau.}
\end{mcq}

\begin{mcq}{Trong không gian $Oxyz$, cho hai đường thẳng $d_1: \begin{cases} x = -2 + 5t \\ y = 1 - t \\ z = 3t \end{cases}$ ($t \in \mathbb{R}$) và $d_2: \frac{x+2}{4} = \frac{y-1}{5} = \frac{z-1}{-6}$. Vị trí tương đối của hai đường thẳng $d_1$ và $d_2$ là:}{0}{1}{$\vec{u_1} = (5; -1; 3), \vec{u_2} = (4; 5; -6)$. Xét hệ phương trình tìm giao điểm ta có nghiệm duy nhất $t = 0, t' = 0 \Rightarrow$ Giao điểm là $(-2; 1; 0)$. Vậy $d_1$ và $d_2$ cắt nhau. Chọn A.}
	\option{Cắt nhau.}
	\option{Chéo nhau.}
	\option{Song song.}
	\option{Trùng nhau.}
\end{mcq}

\begin{mcq}{Trong không gian $Oxyz$, cho hai đường thẳng $d_1: \frac{x}{3} = \frac{y+5}{2} = \frac{z-1}{-3}$ và $d_2: \frac{x-1}{-6} = \frac{y-3}{-4} = \frac{z-1}{6}$. Vị trí tương đối của hai đường thẳng $d_1$ và $d_2$ là:}{1}{1}{$\vec{u_1} = (3; 2; -3), \vec{u_2} = (-6; -4; 6) = -2\vec{u_1}$. Điểm $A(0; -5; 1) \in d_1$ thay vào $d_2$ ta thấy không thỏa mãn. Vậy $d_1 \parallel d_2$. Chọn B.}
	\option{Trùng nhau.}
	\option{Song song.}
	\option{Cắt nhau.}
	\option{Chéo nhau.}
\end{mcq}

\begin{mcq}{Trong không gian $Oxyz$, đường thẳng $d: \begin{cases} x = 1 + t \\ y = 2 - t \\ z = 2t \end{cases}$ ($t \in \mathbb{R}$) vuông góc với đường thẳng nào sau đây?}{1}{1}{$\vec{u_d} = (1; -1; 2)$. Xét đường thẳng $\Delta_2: \frac{x+1}{4} = \frac{y-1}{6} = \frac{z}{1}$ có VTCP $\vec{u_2} = (4; 6; 1)$, ta có $\vec{u_d} \cdot \vec{u_2} = 1(4) + (-1)6 + 2(1) = 0 \Rightarrow d \perp \Delta_2$. Chọn B.}
	\option{$\Delta_1: \frac{x}{2} = \frac{y-2}{4} = \frac{z-1}{6}$.}
	\option{$\Delta_2: \frac{x+1}{4} = \frac{y-1}{6} = \frac{z}{1}$.}
	\option{$\Delta_3: \frac{x}{-1} = \frac{y+2}{-2} = \frac{z}{3}$.}
	\option{$\Delta_4: \frac{x-1}{1} = \frac{y-2}{-1} = \frac{z}{2}$.}
\end{mcq}

\begin{mcq}{Trong không gian $Oxyz$, côsin của góc giữa hai đường thẳng $d: \frac{x-2}{1} = \frac{y+3}{2} = \frac{z-6}{2}$ và $d': \begin{cases} x = 1 - 2t \\ y = 2 - 3t \\ z = 1 - 6t \end{cases}$ ($t \in \mathbb{R}$) bằng}{3}{1}{$\vec{u} = (1; 2; 2), \vec{u'} = (-2; -3; -6)$. $\cos(d, d') = \frac{|1(-2) + 2(-3) + 2(-6)|}{\sqrt{1+4+4}\sqrt{4+9+36}} = \frac{|-20|}{3 \cdot 7} = \frac{20}{21}$. Chọn D.}
	\option{$\frac{4}{21}$.}
	\option{$-\frac{20}{21}$.}
	\option{$\frac{2}{7}$.}
	\option{$\frac{20}{21}$.}
\end{mcq}

\begin{mcq}{Trong không gian $Oxyz$, góc giữa hai mặt phẳng $(\alpha): x + 2y + 2z - 1 = 0$ và $(\beta): 7x - 8y - 15z + 3 = 0$ bằng}{1}{1}{$\vec{n_1} = (1; 2; 2), \vec{n_2} = (7; -8; -15)$. $\cos((\alpha), (\beta)) = \frac{|1(7) + 2(-8) + 2(-15)|}{\sqrt{1+4+4}\sqrt{49+64+225}} = \frac{39}{3 \cdot \sqrt{338}} = \frac{13}{13\sqrt{2}} = \frac{\sqrt{2}}{2} \Rightarrow ((\alpha), (\beta)) = 45^\circ$. Chọn B.}
	\option{$30^\circ$.}
	\option{$45^\circ$.}
	\option{$60^\circ$.}
	\option{$90^\circ$.}
\end{mcq}

\begin{mcq}{Trong không gian $Oxyz$, cho đường thẳng $d: \frac{x}{2} = \frac{y-1}{1} = \frac{z}{1}$ và mặt phẳng $(P): 4x + 3y + 5z - 4 = 0$. Gọi $\alpha$ là góc giữa đường thẳng $d$ và mặt phẳng $(P)$. Khẳng định nào sau đây đúng?}{2}{1}{$\vec{u} = (2; 1; 1), \vec{n} = (4; 3; 5)$. $\sin \alpha = \frac{|2(4) + 1(3) + 1(5)|}{\sqrt{4+1+1}\sqrt{16+9+25}} = \frac{16}{\sqrt{6}\sqrt{50}} = \frac{16}{10\sqrt{3}} = \frac{8\sqrt{3}}{15} \approx 0{,}9238 \Rightarrow \alpha \approx 67^\circ28'$. Chọn C.}
	\option{$\alpha = 30^\circ$.}
	\option{$\alpha = 45^\circ$.}
	\option{$\alpha \approx 67^\circ28'$.}
	\option{$\alpha = 90^\circ$.}
\end{mcq}

\end{quiz}

% =========================================================================
% 2. CÂU TRẮC NGHIỆM ĐÚNG SAI
% =========================================================================

\begin{quiz}{Phần 2. Trắc nghiệm Đúng/Sai}

\begin{truefalse}{Trong không gian $Oxyz$, cho đường thẳng $\Delta: \frac{x+2024}{4} = \frac{y+2025}{-2} = \frac{z+2026}{1}$ và mặt phẳng $(P): x - 2y - 2z + 1 = 0$. Gọi $\alpha$ là góc giữa đường thẳng $\Delta$ và mặt phẳng $(P)$.}{1}{Đường thẳng $\Delta$ có VTCP $\vec{u} = (4; -2; 1)$, mặt phẳng $(P)$ có VTPT $\vec{n} = (1; -2; -2)$. Ta có $\sin \alpha = \frac{|4(1) + (-2)(-2) + 1(-2)|}{\sqrt{16+4+1}\sqrt{1+4+4}} = \frac{6}{\sqrt{21} \cdot 3} = \frac{2}{\sqrt{21}} \approx 0{,}4364 \Rightarrow \alpha \approx 26^\circ$.}
	\statement{false}{Vecto $\vec{u} = (2024; 2025; 2026)$ là một vecto chỉ phương của đường thẳng $\Delta$.}
	\statement{false}{Vecto $\vec{n} = (1; 2; 2)$ là một vecto pháp tuyến của mặt phẳng $(P)$.}
	\statement{true}{Công thức tính sin góc giữa đường thẳng $\Delta$ và mặt phẳng $(P)$ là $\sin \alpha = \frac{|\vec{u} \cdot \vec{n}|}{|\vec{u}| \cdot |\vec{n}|}$.}
	\statement{true}{Góc $\alpha$ giữa đường thẳng $\Delta$ và mặt phẳng $(P)$ xấp xỉ bằng $26^\circ$.}
\end{truefalse}

\begin{truefalse}{Trong không gian $Oxyz$, cho hai mặt phẳng $(P_1): 2x - 3y - 6z + 7 = 0$ và $(P_2): 2x + 2y + z + 8 = 0$. Gọi $\alpha$ là góc giữa hai mặt phẳng $(P_1)$ và $(P_2)$.}{1}{$(P_1)$ có VTPT $\vec{n_1} = (2; -3; -6)$, $(P_2)$ có VTPT $\vec{n_2} = (2; 2; 1)$. Ta có $\cos \alpha = \frac{|2(2) + (-3)(2) + (-6)(1)|}{\sqrt{4+9+36}\sqrt{4+4+1}} = \frac{|4 - 6 - 6|}{7 \cdot 3} = \frac{8}{21} \approx 0{,}3810 \Rightarrow \alpha \approx 68^\circ$.}
	\statement{true}{Vecto $\vec{n_1} = (2; -3; -6)$ là một vecto pháp tuyến của mặt phẳng $(P_1)$.}
	\statement{false}{Vecto $\vec{n_2} = (2; -2; 1)$ là một vecto pháp tuyến của mặt phẳng $(P_2)$.}
	\statement{true}{Công thức tính côsin góc giữa hai mặt phẳng là $\cos \alpha = \frac{|\vec{n_1} \cdot \vec{n_2}|}{|\vec{n_1}| \cdot |\vec{n_2}|}$.}
	\statement{true}{Góc giữa hai mặt phẳng $(P_1)$ và $(P_2)$ xấp xỉ bằng $68^\circ$.}
\end{truefalse}

\begin{truefalse}{Trong không gian $Oxyz$, cho hai đường thẳng $\Delta_1: \frac{x+1}{-1} = \frac{y+4}{2} = \frac{z-5}{-3}$ và $\Delta_2: \frac{x}{2} = \frac{y-3}{-1} = \frac{z+2}{-1}$. Gọi $\alpha$ là góc giữa hai đường thẳng $\Delta_1$ và $\Delta_2$.}{1}{$\Delta_1$ qua $M_1(-1; -4; 5)$ có VTCP $\vec{u_1} = (-1; 2; -3)$. $\Delta_2$ qua $M_2(0; 3; -2)$ có VTCP $\vec{u_2} = (2; -1; -1)$. Ta có $\vec{M_1M_2} \cdot [\vec{u_1}, \vec{u_2}] = 1(-5) + 7(-7) + (-7)(-3) = -33 \ne 0 \Rightarrow \Delta_1, \Delta_2$ chéo nhau. $\cos \alpha = \frac{|-2 - 2 + 3|}{\sqrt{14}\sqrt{6}} = \frac{1}{\sqrt{84}} \approx 0{,}1091 \Rightarrow \alpha \approx 84^\circ$.}
	\statement{false}{Đường thẳng $\Delta_1$ có một vecto chỉ phương là $\vec{n_1} = (-1; -4; 5)$.}
	\statement{true}{Đường thẳng $\Delta_2$ đi qua điểm $A(0; 3; -2)$.}
	\statement{true}{Hai đường thẳng $\Delta_1$ và $\Delta_2$ chéo nhau.}
	\statement{true}{Góc $\alpha$ giữa hai đường thẳng $\Delta_1$ và $\Delta_2$ xấp xỉ bằng $84^\circ$.}
\end{truefalse}

\begin{truefalse}{Trong không gian $Oxyz$, cho điểm $A(1; 2; 3)$ và đường thẳng $d: \begin{cases} x = 3 + t \\ y = 5 - 2t \\ z = 7 - 2t \end{cases}$ ($t \in \mathbb{R}$).}{1}{Thay $A(1; 2; 3)$ vào $d$ ta được $t = -2$ thỏa mãn. VTCP của $d$ là $\vec{u} = (1; -2; -2)$. Đường thẳng $d'$ có VTCP $\vec{u'} = (2; 2; -1)$, ta có $\vec{u} \cdot \vec{u'} = 2 - 4 + 2 = 0 \Rightarrow d \perp d'$. Đường thẳng $\Delta$ qua $(4; 3; 5) \in d$ và có VTCP $(2; -4; -4) = 2\vec{u} \Rightarrow \Delta \equiv d$.}
	\statement{true}{Điểm $A(1; 2; 3)$ thuộc đường thẳng $d$.}
	\statement{true}{Đường thẳng đi qua $A$ và song song với $d$ có phương trình chính tắc là $\frac{x-1}{1} = \frac{y-2}{-2} = \frac{z-3}{-2}$.}
	\statement{true}{Đường thẳng $d$ vuông góc với đường thẳng $d': \frac{x+4}{2} = \frac{y+6}{2} = \frac{z-10}{-1}$.}
	\statement{true}{Đường thẳng $d$ và đường thẳng $\Delta: \frac{x-4}{2} = \frac{y-3}{-4} = \frac{z-5}{-4}$ trùng nhau.}
\end{truefalse}

\end{quiz}

% =========================================================================
% 3. CÂU TRẮC NGHIỆM TRẢ LỜI NGẮN
% =========================================================================

\begin{quiz}{Phần 3. Trắc nghiệm trả lời ngắn}

\shortanswer{Trong không gian $Oxyz$, cho đường thẳng $d$ đi qua điểm $N(-2; 3; 1)$ và nhận $\vec{a} = (3; -4; 5)$ làm vecto chỉ phương. Gọi $A$ là điểm thuộc $d$ có hoành độ bằng $4$. Cao độ của điểm $A$ bằng bao nhiêu?}{11}{1}{Phương trình tham số của $d: \begin{cases} x = -2 + 3t \\ y = 3 - 4t \\ z = 1 + 5t \end{cases}$. Điểm $A \in d$ có hoành độ $x_A = 4 \Leftrightarrow -2 + 3t = 4 \Leftrightarrow t = 2$. Suy ra cao độ của điểm $A$ là $z_A = 1 + 5(2) = 11$.}

\shortanswer{Trong không gian $Oxyz$, cho hình chóp $S.ABCD$ có đáy $ABCD$ là hình bình hành. Biết tọa độ các đỉnh $S(3; -2; 4), A(3; 4; 5), B(8; 8; 6), C(7; 6; 3)$. Góc giữa đường thẳng $d$ chứa cạnh bên $SB$ và đường thẳng $d'$ chứa cạnh đáy $AD$ bằng bao nhiêu độ (làm tròn kết quả đến hàng đơn vị của độ)?}{43}{1}{Ta có $\vec{SB} = (5; 10; 2)$. Vì $ABCD$ là hình bình hành nên $\vec{AD} = \vec{BC} = (-1; -2; -3)$. Côsin góc giữa hai đường thẳng $SB$ và $AD$ là: $\cos(SB, AD) = \frac{|5(-1) + 10(-2) + 2(-3)|}{\sqrt{25+100+4}\sqrt{1+4+9}} = \frac{31}{\sqrt{129}\sqrt{14}} = \frac{31}{\sqrt{1806}} \approx 0{,}7295 \Rightarrow (SB, AD) \approx 43^\circ$.}

\shortanswer{Trong không gian $Oxyz$, cho đường thẳng $d: \begin{cases} x = 2 - 2t \\ y = 3 + 3t \\ z = 6 + 4t \end{cases}$ ($t \in \mathbb{R}$). Gọi điểm $A(x_0; y_0; z_0)$ thuộc đường thẳng $d$ cách gốc tọa độ $O$ một khoảng bằng $7$, biết $z_0$ là một số nguyên âm. Tổng $x_0 + y_0 + z_0$ bằng bao nhiêu?}{1}{1}{Điểm $A \in d \Rightarrow A(2 - 2t; 3 + 3t; 6 + 4t)$. $OA = 7 \Leftrightarrow (2 - 2t)^2 + (3 + 3t)^2 + (6 + 4t)^2 = 49 \Leftrightarrow 29t^2 + 58t = 0 \Leftrightarrow \begin{cases} t = 0 \\ t = -2 \end{cases}$. Vì $z_0$ là số nguyên âm nên ta chọn $t = -2 \Rightarrow A(6; -3; -2)$ có $z_0 = -2 < 0$. Khi đó $x_0 + y_0 + z_0 = 6 + (-3) + (-2) = 1$.}

\shortanswer{\image{0.6}{Mô hình đường trượt zipline trong không gian $Oxyz$}{images_ChuDe04_PhuongTrinhDuongThang/lt_14.png} Một đường trượt zipline trong một khu du lịch sinh thái nối từ điểm xuất phát $A(3; 2{,}5; 15)$ đến điểm kết thúc $B(21; 27{,}5; 10)$ (đơn vị trên các trục tính bằng mét). Một du khách đang ở vị trí điểm $M(x_0; y_0; 12)$ trên đường trượt ở độ cao $12\text{ m}$. Giá trị của tích $x_0 y_0$ bằng bao nhiêu?}{241,5}{1}{Đường thẳng $AB$ đi qua $A(3; 2{,}5; 15)$ và có VTCP $\vec{AB} = (18; 25; -5)$ nên có phương trình: $\begin{cases} x = 3 + 18t \\ y = 2{,}5 + 25t \\ z = 15 - 5t \end{cases}$. Tại vị trí điểm $M$ có $z = 12 \Leftrightarrow 15 - 5t = 12 \Leftrightarrow t = 0{,}6$. Khi đó $x_0 = 3 + 18(0{,}6) = 13{,}8$ và $y_0 = 2{,}5 + 25(0{,}6) = 17{,}5$. Tích $x_0 y_0 = 13{,}8 \times 17{,}5 = 241{,}5$.}

\shortanswer{\image{0.5}{Kim tự tháp Kheops}{images_ChuDe04_PhuongTrinhDuongThang/lt_15.png} Kim tự tháp Kheops (Ai Cập) có dạng hình chóp đều $S.ABCD$ với đáy $ABCD$ là hình vuông cạnh $230\text{ m}$, các cạnh bên bằng nhau và dài $219\text{ m}$. Góc giữa hai mặt bên kề nhau $(SAB)$ và $(SBC)$ của kim tự tháp bằng bao nhiêu độ (làm tròn kết quả đến hàng đơn vị của độ)?}{68}{1}{Chọn hệ trục tọa độ $Oxyz$ với gốc $O$ là tâm đáy hình vuông $ABCD$, trục $Oz$ trùng với trục đối xứng $SO$. Tọa độ các đỉnh: $A(-115; -115; 0), B(115; -115; 0), C(115; 115; 0), D(-115; 115; 0)$ và $S(0; 0; h)$ với $h = \sqrt{219^2 - 115^2 - 115^2} \approx 146{,}67\text{ m}$. Mặt phẳng $(SAB)$ có VTPT $\vec{n_1} = (0; -h; 115)$, mặt phẳng $(SBC)$ có VTPT $\vec{n_2} = (h; 0; 115)$. Ta có $\cos((SAB), (SBC)) = \frac{115^2}{h^2 + 115^2} = \frac{115^2}{219^2 - 115^2} \approx 0{,}3807 \Rightarrow ((SAB), (SBC)) \approx 68^\circ$.}

\shortanswer{\image{0.6}{Mô hình khuôn đúc bánh mì trong không gian $Oxyz$}{images_ChuDe04_PhuongTrinhDuongThang/lt_16.png} Một chiếc khuôn đúc bánh mì được mô hình hóa trong không gian $Oxyz$. Đáy khuôn là hình chữ nhật $SPQR$ với $S(0; 0; 0), P(8; 0; 0), Q(8; 18; 0)$. Miệng khuôn mở rộng với các cạnh bên nghiêng đều, đỉnh mép tương ứng là $T(-1; -1; 7)$. Gọi $m$ là góc giữa cạnh bên $ST$ và cạnh đáy $SP$, $n$ là góc giữa cạnh bên $ST$ và mặt đáy $(Oxy)$, $p$ là góc giữa mặt bên chứa cạnh $ST, SP$ và mặt đáy $(Oxy)$. Tổng $m + n + p$ bằng bao nhiêu độ (làm tròn kết quả đến hàng đơn vị của độ)?}{242}{1}{Ta có $\vec{ST} = (-1; -1; 7), \vec{SP} = (8; 0; 0)$. 
- Góc $m$ giữa cạnh bên $ST$ và cạnh đáy $SP$: $\cos m = \frac{|-8|}{\sqrt{51} \cdot 8} = \frac{1}{\sqrt{51}} \Rightarrow m \approx 81{,}95^\circ$.
- Góc $n$ giữa cạnh bên $ST$ và mặt đáy $(Oxy)$: $\sin n = \frac{7}{\sqrt{51}} \Rightarrow n \approx 78{,}58^\circ$.
- Góc $p$ giữa mặt bên $(STP)$ và mặt đáy $(Oxy)$: VTPT mặt bên là $\vec{n_{(STP)}} = [\vec{SP}, \vec{ST}] = (0; -56; -8) = -8(0; 7; 1) \Rightarrow \cos p = \frac{1}{\sqrt{50}} \Rightarrow p \approx 81{,}87^\circ$.
Tổng $m + n + p \approx 81{,}95^\circ + 78{,}58^\circ + 81{,}87^\circ = 242{,}4^\circ \approx 242^\circ$.}

\end{quiz}

\end{document}
